\chapter{Ortogonalność i rzuty}
\label{ch:lin13-ortogonalnosc-rzuty}
\index{rzut ortogonalny@Rzut ortogonalny}
\index{dopełnienie ortogonalne@Dopełnienie ortogonalne}
\index{przybliżenie najlepsze@Przybliżenie najlepsze}
\index{równania normalne@Równania normalne}
\index{operator sprzężony@Operator sprzężony}

\noindent\textbf{Cel rozdziału.}
W Rozdziale~\ref{ch:lin9-formy-kwadratowe}
poznaliśmy dopełnienia ortogonalne dla rzeczywistych
form symetrycznych i udowodniliśmy istnienie rzutów
na podprzestrzenie w przestrzeniach euklidesowych.
W Rozdziale~\ref{ch:lin12-euklidesowe-unitarne}
rozszerzyliśmy iloczyny skalarne na ciało zespolone,
udowodniliśmy ortogonalizację i reprezentację Riesza.
Nie powtarzamy tych dowodów.
Celem obecnego rozdziału jest opis projekcji za pomocą
\emph{operatorów i macierzy}, rozwiązywanie problemu
najlepszego przybliżenia, równania normalne i
metoda najmniejszych kwadratów.
Wprowadzimy też operator sprzężony jako narzędzie
porządkujące wszystkie te konstrukcje.

Przez \(K\) oznaczamy \(\mathbb R\) lub \(\mathbb C\).
Przestrzenie \((V,h_V)\), \((W,h_W)\)
mają dodatnio określone iloczyny skalarne,
liniowe w \emph{drugim} argumencie:
\(h_V(au,v)=\overline a\,h_V(u,v)\),
\(h_V(u,av)=a h_V(u,v)\).
Wymiar skończony zakładamy tam, gdzie jest potrzebny,
a nie bez uzasadnienia dla całego rozdziału.
W bazach ortonormalnych zapisujemy macierze
przez \(A^*\), czyli sprzężenie transpozycję.

\section{Dopełnienia ortogonalne nad \(\mathbb R\) i \(\mathbb C\)}
\label{sec:lin13-dopelnienia}

\begin{definicja}[Dopełnienie ortogonalne]
\label{def:lin13-dopelnienie}
Niech \((V,h)\) będzie przestrzenią euklidesową
lub unitarną, a \(U\subseteq V\) podprzestrzenią.
Definiujemy
\[
 U^\perp=\{v\in V:\ h(u,v)=0
                     \text{ dla wszystkich }u\in U\}.
\]
Z hermitowskości wynika równoważność
\(h(u,v)=0\iff h(v,u)=0\);
kolejność argumentów w definicji
nie zmienia więc podprzestrzeni.
\end{definicja}

\begin{stwierdzenie}[Algebra dopełnień ortogonalnych]
\label{prop:lin13-dopelnienia}
Niech \((V,h)\) będzie przestrzenią
euklidesową lub unitarną.
Dla dowolnych podprzestrzeni
\(U,U_1,U_2\subseteq V\) mamy
\[
 \begin{gathered}
 U^\perp\text{ jest podprzestrzenią},\qquad
 U\cap U^\perp=\{0\},\qquad
 U\subseteq (U^\perp)^\perp,\\
 U_1\subseteq U_2\Longrightarrow
     U_2^\perp\subseteq U_1^\perp,\\
 (U_1+U_2)^\perp=U_1^\perp\cap U_2^\perp.
 \end{gathered}
\]
Nie zakładamy skończonego wymiaru.
\end{stwierdzenie}
\begin{proof}
Dla każdego \(u\in U\) odwzorowanie
\(v\mapsto h(u,v)\) jest liniowe,
więc \(U^\perp\) jest przecięciem
jąder odwzorowań liniowych.
Jeśli \(v\in U\cap U^\perp\),
to \(h(v,v)=0\); dodatnia określoność
daje \(v=0\).
Jeżeli \(u\in U\) i \(v\in U^\perp\),
to z hermitowskości
\(h(v,u)=\overline{h(u,v)}=0\),
skąd \(u\in(U^\perp)^\perp\).
Przy \(U_1\subseteq U_2\)
każdy wektor ortogonalny do \(U_2\)
jest ortogonalny do \(U_1\).

Jeśli \(v\in(U_1+U_2)^\perp\),
jest ortogonalny do każdego
elementu obu składników,
ponieważ \(U_i\subseteq U_1+U_2\).
Odwrotnie, dla
\(v\in U_1^\perp\cap U_2^\perp\)
i \(u_1\in U_1,u_2\in U_2\)
liniowość w pierwszym argumencie
z odpowiednim sprzężeniem daje
\(h(u_1+u_2,v)=0\).
Stąd obie strony ostatniej
równości są równe.
\end{proof}

\begin{twierdzenie}[Rozkład względem podprzestrzeni skończonego wymiaru]
\label{thm:lin13-rozklad}
Niech \((V,h)\) będzie przestrzenią
euklidesową lub unitarną,
niekoniecznie skończenie wymiarową.
Niech \(U\subseteq V\) będzie
podprzestrzenią \emph{skończonego wymiaru}.
Wtedy istnieje jednoznaczny rozkład
\[
 v=u+w,\qquad u\in U,\ w\in U^\perp,
 \quad\text{dla każdego }v\in V,
\]
czyli \(V=U\oplus U^\perp\).
Ponadto
\[
 (U^\perp)^\perp=U.
\]
W szczególności jeśli \(\dim V=n<\infty\),
to dla dowolnego \(U\subseteq V\)
zachodzi
\(\dim U^\perp=n-\dim U\).
\end{twierdzenie}
\begin{proof}
Wybierzmy ortonormalną bazę
\((e_1,\ldots,e_r)\) przestrzeni \(U\).
Jej istnienie w przypadku zespolonym
zapewnia
Twierdzenie~\ref{thm:lin12-gram-schmidt},
a w rzeczywistym
Twierdzenie~\ref{thm:lin9-gram-schmidt}.
Dla \(v\in V\) połóżmy
\[
 u=\sum_{j=1}^r h(e_j,v)e_j,
 \qquad w=v-u.
\]
Wtedy \(u\in U\), a dla każdego \(i\)
\[
 h(e_i,w)
 =h(e_i,v)-\sum_j h(e_j,v)h(e_i,e_j)=0.
\]
Ponieważ \(e_i\) rozpinają \(U\),
otrzymujemy \(w\in U^\perp\).
Jeśli \(v=u_1+w_1=u_2+w_2\)
są dwa takie rozkłady, to
\(u_1-u_2=w_2-w_1\)
należy do \(U\cap U^\perp=\{0\}\),
więc \(u_1=u_2\) i \(w_1=w_2\).

Dowiedliśmy już
\(U\subseteq(U^\perp)^\perp\).
Jeśli \(z\in(U^\perp)^\perp\),
piszemy \(z=u+w\) jak wyżej.
Wektor \(w\in U^\perp\), zatem
\(h(w,z)=0\); jednocześnie
\(h(w,u)=0\). Stąd
\(0=h(w,z)=h(w,u+w)=h(w,w)\),
więc \(w=0\), czyli \(z=u\in U\).
To dowodzi równości podprzestrzeni.
Dla skończenie wymiarowego \(V\)
rozmiar sumy prostej daje
\(\dim V=\dim U+\dim U^\perp\).
Przy \(U=\{0\}\) podane sumy są puste,
\(U^\perp=V\), i wszystkie wnioski
nadal zachodzą.
\end{proof}

\begin{uwaga}[Zakres twierdzenia o rozkładzie]
\label{rem:lin13-nieskonczony}
Skończony wymiar podprzestrzeni \(U\)
jest tu istotnym, wystarczającym założeniem;
nie wyciągamy z dowodu wniosku
\(V=U\oplus U^\perp\) dla \emph{dowolnej}
podprzestrzeni nieskończenie wymiarowej
w przestrzeni z iloczynem skalarnym.
Dla przestrzeni Hilberta istnieje
szersze twierdzenie dotyczące
podprzestrzeni domkniętych,
lecz wymaga ono topologii i zupełności.
Obecny dowód używa wyłącznie
skończonej bazy ortonormalnej \(U\).
\end{uwaga}

\section{Operator rzutu ortogonalnego}
\label{sec:lin13-projekcja}

\begin{definicja}[Rzut na podprzestrzeń]
\label{def:lin13-rzut}
Niech \(U\subseteq V\) będzie
skończenie wymiarową podprzestrzenią
przestrzeni euklidesowej lub unitarnej.
\emph{Rzutem ortogonalnym na \(U\)}
nazywamy odwzorowanie
\(P_U\colon V\to V\) przyporządkowujące
wektorowi \(v\) składnik \(u\in U\)
jego rozkładu \(v=u+w\),
\(w\in U^\perp\), z
Twierdzenia~\ref{thm:lin13-rozklad}.
\end{definicja}

\begin{twierdzenie}[Wzory na rzut i jego jądro]
\label{thm:lin13-rzut-wlasnosci}
Niech \((V,h)\) będzie przestrzenią
euklidesową lub unitarną, a
\(U\subseteq V\) podprzestrzenią
skończonego wymiaru.
Wtedy \(P_U\) jest liniowe i
\[
 \begin{gathered}
 \operatorname{im}P_U=U,\qquad
 \ker P_U=U^\perp,\qquad
 P_U^2=P_U,\\
 h(P_Uv,w)=h(v,P_Uw)
 \quad(v,w\in V).
 \end{gathered}
\]
Jeśli \((e_1,\ldots,e_r)\)
jest bazą ortonormalną \(U\), to
\[
 P_Uv=\sum_{j=1}^r h(e_j,v)e_j.
\]
Rzut \(P_U\) nie zależy od
wyboru bazy ortonormalnej \(U\).
Rzut \(I_V-P_U\) ma obraz \(U^\perp\)
i jądro \(U\), niezależnie
od wymiaru \(U^\perp\).
\end{twierdzenie}
\begin{proof}
Liniowość wynika wprost
z liniowości współczynników
\(h(e_j,v)\) w drugim argumencie.
Z
Twierdzenia~\ref{thm:lin13-rozklad}
podany wzór wybiera dokładnie
składnik należący do \(U\).
Jeśli \(u\in U\), to rozkład
\(u=u+0\) daje \(P_Uu=u\),
więc obraz to \(U\) i
\(P_U^2v=P_Uv\).
Równocześnie \(P_Uv=0\)
dokładnie wtedy, gdy \(v\in U^\perp\),
zatem \(\ker P_U=U^\perp\).

Dla \(v=u_v+w_v\) i
\(w=u_w+w_w\) z \(u_v,u_w\in U\)
i \(w_v,w_w\in U^\perp\) mamy
\[
 h(P_Uv,w)=h(u_v,u_w)
          =h(v,P_Uw).
\]
Jednoznaczność rozkładu
gwarantuje niezależność \(P_U\)
od użytej bazy.
Odwzorowanie \(I_V-P_U\)
przyporządkowuje \(v\)
składnik \(w_v\), więc ma
obraz \(U^\perp\), jądro \(U\)
i jest idempotentne.
\end{proof}

\begin{twierdzenie}[Najlepsze przybliżenie]
\label{thm:lin13-najlepsze}
Niech \((V,h)\) będzie
przestrzenią euklidesową lub unitarną,
a \(U\subseteq V\) podprzestrzenią
skończonego wymiaru.
Dla każdego \(v\in V\)
i \(u\in U\) zachodzi
\[
 \|v-u\|^2
 =\|v-P_Uv\|^2+\|P_Uv-u\|^2.
\]
Zatem \(P_Uv\) jest jedynym
punktem podprzestrzeni \(U\)
minimalizującym odległość do \(v\).
Dokładniej,
\[
 \operatorname{dist}(v,U)
 :=\inf_{u\in U}\|v-u\|
 =\|v-P_Uv\|.
\]
\end{twierdzenie}
\begin{proof}
Z własności rzutu
\(v-P_Uv\in U^\perp\).
Dla \(u\in U\) mamy też
\(P_Uv-u\in U\), więc te
dwa wektory są ortogonalne.
Z
Wniosku~\ref{cor:lin12-pitagoras}
otrzymujemy wskazaną równość.
Ostatni składnik jest nieujemny
i znika dokładnie wtedy,
gdy \(u=P_Uv\).
Wobec tego minimum jest
osiągane właśnie w \(P_Uv\)
i nigdzie indziej.
Równość dla odległości
wynika z wyciągnięcia
pierwiastków z równości norm.
\end{proof}

\begin{stwierdzenie}[Kontrakcja i optymalny rozkład normy]
\label{prop:lin13-kontrakcja}
Niech \(P_U\) będzie rzutem
ortogonalnym na skończenie
wymiarową podprzestrzeń
przestrzeni euklidesowej
lub unitarnej \(V\).
Dla każdego \(v\in V\) zachodzi
\[
 \|P_Uv\|\le\|v\|,\qquad
 \|(I_V-P_U)v\|\le\|v\|,
 \qquad
 \|v\|^2=\|P_Uv\|^2+
             \|(I_V-P_U)v\|^2.
\]
Jeżeli \(U\ne\{0\}\), to
\(\sup_{v\ne0}\|P_Uv\|/\|v\|=1\).
Jeśli \(U=\{0\}\), to \(P_U=0\).
\end{stwierdzenie}
\begin{proof}
Z rozkładu ortogonalnego
\(v=P_Uv+(I_V-P_U)v\)
wynika wzór Pitagorasa.
Oba kwadraty norm są
nieujemne, co daje obie
nierówności. Jeśli
\(U\ne\{0\}\), wybieramy
\(0\ne u\in U\). Ponieważ
\(P_Uu=u\), iloraz norm
jest równy \(1\), a
nie może przekraczać \(1\).
Dla \(U=\{0\}\) obraz rzutu
jest zerowy, więc rzut
jest odwzorowaniem zerowym.
\end{proof}

\begin{przyklad}[Rzut na prostą i hiperpłaszczyznę]
\label{prz:lin13-prosta}
Niech \(V=K^n\) ze standardowym
iloczynem skalarnym i niech \(a\in K^n\)
będzie niezerowy. Rzut na prostą
\(U=\operatorname{span}\{a\}\) ma postać
\[
 P_Uv=\frac{a^*v}{a^*a}\,a.
\]
Wynika to z wyboru
\(e=a/\sqrt{a^*a}\)
i wzoru na rzut w bazie
ortonormalnej \((e)\).
Dla hiperpłaszczyzny
\(H=a^\perp=\{v:a^*v=0\}\)
otrzymujemy
\[
 P_Hv=v-\frac{a^*v}{a^*a}\,a.
\]
Są to wzory poprawne także
nad \(\mathbb C\), ponieważ
\(a^*v\), a nie \(a^{\mathsf T}v\),
jest współczynnikiem rzutu.
\end{przyklad}

\begin{figure}[!htbp]
\centering
\begin{tikzpicture}[>=Latex,scale=.9]
 \draw[->,gray] (-0.35,0)--(5.7,0)
   node[right] {\(e_1\)};
 \draw[->,gray] (0,-0.35)--(0,3.6)
   node[above] {\(e_2\)};
 \draw[manifoldblue,very thick] (-.3,-.15)--(5.25,2.625);
 \node[manifoldblue,above right] at (4.7,2.35) {\(U\)};
 \coordinate (O) at (0,0);
 \coordinate (P) at (2.8,1.4);
 \coordinate (V) at (2.05,2.9);
 \draw[->,very thick,manifoldblue] (O)--(P)
    node[pos=.65,below right] {\(P_Uv\)};
 \draw[->,very thick,manifoldred] (O)--(V)
    node[pos=.80,left] {\(v\)};
 \draw[densely dashed,very thick,manifoldgreen] (P)--(V)
    node[pos=.64,right] {\(v-P_Uv\)};
 \draw[manifoldblue!70!black] ($(P)!0.14!(O)$)
   -- ++(-.145,.29) -- ++(.28,.14);
 \fill (P) circle (1.8pt);
\end{tikzpicture}
\caption{Rzut ortogonalny \(P_Uv\) i reszta
\(v-P_Uv\), prostopadła do podprzestrzeni \(U\).
Wektor \(P_Uv\) jest najlepszym
przybliżeniem \(v\) w \(U\).}
\label{fig:lin13-rzut-geometryczny}
\end{figure}


\section{Operator sprzężony i rzutowanie w skończonym wymiarze}
\label{sec:lin13-operator-sprzezony}

\begin{twierdzenie}[Istnienie operatora sprzężonego]
\label{thm:lin13-adjoint}
Niech \((V,h_V)\), \((W,h_W)\)
będą skończenie wymiarowymi
przestrzeniami euklidesowymi
lub unitarnymi nad tym samym ciałem \(K\).
Dla każdego odwzorowania liniowego
\(A\colon V\to W\)
istnieje dokładnie jedno liniowe
\(A^*\colon W\to V\), zwane
\emph{operatorem sprzężonym},
dla którego
\[
 h_W(Av,w)=h_V(v,A^*w)
 \qquad(v\in V,\ w\in W).
\]
W bazach ortonormalnych \(V,W\)
macierz operatora \(A^*\) jest
sprzężoną transpozycją macierzy \(A\).
\end{twierdzenie}
\begin{proof}
Ustalmy \(w\in W\).
Odwzorowanie
\[
 F_w\colon V\to K,\qquad
 F_w(v)=h_W(w,Av)
\]
jest liniowe względem \(v\),
ponieważ \(A\) jest liniowe,
a \(h_W\) jest liniowe
w drugim argumencie.
Z
Twierdzenia~\ref{thm:lin12-riesz}
istnieje dokładnie jeden
wektor \(a_w\in V\) taki, że
\[
 h_V(a_w,v)=h_W(w,Av)
 \qquad(v\in V).
\]
Po sprzężeniu obu stron
otrzymujemy
\[
 h_V(v,a_w)=h_W(Av,w).
\]
Określamy \(A^*w=a_w\).
Pozostaje sprawdzić liniowość.
Dla \(\lambda,\mu\in K\)
i \(w_1,w_2\in W\)
mamy dla każdego \(v\in V\)
\[
 \begin{aligned}
 h_V(v,A^*(\lambda w_1+\mu w_2))
 &=h_W(Av,\lambda w_1+\mu w_2)\\
 &=\lambda h_W(Av,w_1)+
   \mu h_W(Av,w_2)\\
 &=h_V(v,\lambda A^*w_1+\mu A^*w_2).
 \end{aligned}
\]
Podstawiając za \(v\)
różnicę dwóch wektorów
w drugim argumencie,
z dodatniej określoności
uzyskujemy ich równość.
Stąd \(A^*\) jest liniowe.
Ten sam argument dowodzi
jego jednoznaczności.

Niech teraz
\((e_1,\ldots,e_n)\),
\((f_1,\ldots,f_m)\)
będą bazami ortonormalnymi
\(V,W\), a \(M=[A]\).
Wtedy \([Av]=MX\),
\([w]=Y\) i
\[
 h_W(Av,w)=(MX)^*Y=X^*M^*Y.
\]
Z drugiej strony
\(h_V(v,A^*w)=X^*[A^*]Y\).
Równość dla wszystkich
\(X,Y\) daje \([A^*]=M^*\).
Przy wymiarze \(0\)
odwzorowania między odpowiednimi
przestrzeniami są jednoznaczne,
więc teza również obowiązuje.
\end{proof}

\begin{stwierdzenie}[Rachunek operatorów sprzężonych]
\label{prop:lin13-adjoint-kalkul}
Niech \(V,W,Z\) będą
skończenie wymiarowymi
przestrzeniami euklidesowymi
lub unitarnymi nad \(K\),
a \(A,B\colon V\to W\),
\(C\colon W\to Z\) odwzorowaniami liniowymi.
Dla dowolnych \(\lambda,\mu\in K\)
zachodzą
\[
 \begin{gathered}
 (\lambda A+\mu B)^*
   =\overline\lambda A^*+
    \overline\mu B^*,\\
 (C\circ A)^*=A^*\circ C^*,\qquad
 (A^*)^*=A,\qquad
 (I_V)^*=I_V.
 \end{gathered}
\]
Ponadto dla każdego \(v\in V\)
\[
 h_V(v,A^*Av)=h_W(Av,Av)\ge0.
\]
\end{stwierdzenie}
\begin{proof}
Dla pierwszego wzoru
stosujemy sprzężenie liniowość
pierwszego argumentu:
\[
 \begin{aligned}
 h_W((\lambda A+\mu B)v,w)
 &=\overline\lambda\,h_W(Av,w)
   +\overline\mu\,h_W(Bv,w)\\
 &=h_V(v,(\overline\lambda A^*+
                   \overline\mu B^*)w).
 \end{aligned}
\]
Jednoznaczność sprzężenia daje
pierwszą tożsamość.
Dla złożenia mamy
\[
 h_Z(CAv,z)
 =h_W(Av,C^*z)
 =h_V(v,A^*C^*z).
\]
Ponownie stosujemy jednoznaczność.
Z hermitowskości
\[
 h_W(w,Av)=\overline{h_W(Av,w)}
 =\overline{h_V(v,A^*w)}
 =h_V(A^*w,v),
\]
stąd \((A^*)^*=A\).
Wzór dla identyczności wynika
z \(h_V(v,I_Vw)=h_V(I_Vv,w)\).
W końcu
\(h_V(v,A^*Av)=h_W(Av,Av)\)
jest rzeczywisty i nieujemny
przez dodatnią określoność.
\end{proof}

\begin{twierdzenie}[Jądro i obraz operatora sprzężonego]
\label{thm:lin13-kernel-image}
Niech \(V,W\) będą skończenie
wymiarowymi przestrzeniami
euklidesowymi albo unitarnymi
nad \(K\), a \(A\colon V\to W\)
odwzorowaniem liniowym.
Wtedy
\[
 \begin{gathered}
 (\operatorname{im}A)^\perp=\ker A^*,
 \qquad
 (\operatorname{im}A^*)^\perp=\ker A,\\
 \operatorname{im}A=(\ker A^*)^\perp,
 \qquad
 \operatorname{im}A^*=(\ker A)^\perp,\\
 \ker(A^*A)=\ker A,\qquad
 \operatorname{rk}(A^*A)=\operatorname{rk}A.
 \end{gathered}
\]
Dopełnienia w pierwszej kolumnie
są brane odpowiednio w \(W,V\).
\end{twierdzenie}
\begin{proof}
Dla \(w\in W\) mamy
\[
 \begin{aligned}
 w\in(\operatorname{im}A)^\perp
 &\Longleftrightarrow
   h_W(Av,w)=0\quad(v\in V)\\
 &\Longleftrightarrow
   h_V(v,A^*w)=0\quad(v\in V)\\
 &\Longleftrightarrow A^*w=0.
 \end{aligned}
\]
Ostatnia równoważność wynika
z dodatniej określoności,
po podstawieniu \(v=A^*w\).
Stosując ten sam argument
do \(A^*\) i korzystając z
\((A^*)^*=A\), otrzymujemy
\((\operatorname{im}A^*)^\perp=\ker A\).
Oba obrazy są skończenie
wymiarowymi podprzestrzeniami.
Dlatego przez dwukrotne
dopełnienie ortogonalne
z
Twierdzenia~\ref{thm:lin13-rozklad}
otrzymujemy dwie następne równości.

Jeśli \(Av=0\), to
\(A^*Av=0\), więc
\(\ker A\subseteq\ker(A^*A)\).
W drugą stronę, jeśli
\(A^*Av=0\), to
\[
 0=h_V(v,A^*Av)=h_W(Av,Av),
\]
stąd \(Av=0\).
Jądra są równe; twierdzenie
o rzędzie i jądrze
na \(V\) daje równość rzędów.
\end{proof}

\begin{twierdzenie}[Charakterystyka rzutu ortogonalnego]
\label{thm:lin13-projection-characterization}
Niech \((V,h)\) będzie skończenie
wymiarową przestrzenią euklidesową
lub unitarną i \(P\colon V\to V\)
odwzorowaniem liniowym.
Wówczas \(P\) jest rzutem
ortogonalnym na \(\operatorname{im}P\)
dokładnie wtedy, gdy
\[
 P^2=P\quad\text{oraz}\quad P^*=P.
\]
Jeżeli te warunki są spełnione,
to w dowolnej bazie ortonormalnej
macierz \(P\) spełnia
\(M^2=M=M^*\).
\end{twierdzenie}
\begin{proof}
Dla rzutu ortogonalnego
\(P=P_U\) idempotencję
oraz tożsamość
\(h(Pv,w)=h(v,Pw)\)
udowodniono w
Twierdzeniu~\ref{thm:lin13-rzut-wlasnosci}.
Druga jest dokładnie
definicją \(P^*=P\).

Załóżmy odwrotnie,
że \(P^2=P\) i \(P^*=P\).
Z idempotencji wynika
\[
 V=\operatorname{im}P\oplus\ker P.
\]
Istotnie dla każdego \(v\)
mamy \(v=Pv+(v-Pv)\), gdzie
\(Pv\in\operatorname{im}P\),
\(P(v-Pv)=Pv-P^2v=0\);
jeśli \(u=Px\) należy
zarazem do jądra, to
\(u=Pu=0\).
Dla \(u=Px\in\operatorname{im}P\)
i \(w\in\ker P\) zachodzi
\[
 h(u,w)=h(Px,w)=h(x,P^*w)
       =h(x,Pw)=0.
\]
Zatem obie składowe
są ortogonalne, a z
jednoznaczności rozkładu
\(P\) musi być rzutem
ortogonalnym na swój obraz.
Równości macierzowe wynikają
z własności macierzy złożenia
i sprzężenia w bazie
ortonormalnej.
\end{proof}

\begin{przyklad}[Rzut idempotentny, który nie jest ortogonalny]
\label{prz:lin13-oblique}
Na \(\mathbb R^2\) rozważmy
\[
 P=\begin{pmatrix}1&1\\0&0\end{pmatrix}.
\]
Bezpośrednie mnożenie daje
\(P^2=P\), ale
\(P^{\mathsf T}
=\begin{pmatrix}1&0\\1&0\end{pmatrix}
\ne P\).
Obraz \(P\) to
\(\operatorname{span}\{e_1\}\),
a jądro to
\(\operatorname{span}\{(-1,1)\}\);
obie proste nie są prostopadłe.
W szczególności
\(P(0,1)=(1,0)\),
podczas gdy rzut ortogonalny
na \(\operatorname{span}\{e_1\}\)
wysyła \((0,1)\) na \((0,0)\).
Warunek \(P^2=P\) sam nie
wystarcza do ortogonalności.
\end{przyklad}

\section{Macierz rzutu przy dowolnej bazie podprzestrzeni}
\label{sec:lin13-matrix-projection}

\begin{twierdzenie}[Wzór z macierzą Grama]
\label{thm:lin13-projection-gram}
Niech \(n,r\in\mathbb N_0\)
oraz niech \(B\in\operatorname{Mat}_{n\times r}(K)\)
ma liniowo niezależne kolumny.
Niech \(U=\operatorname{im}B\subseteq K^n\),
gdzie \(K^n\) ma standardowy
iloczyn skalarny.
Wtedy macierz Grama
\(G=B^*B\in\operatorname{Mat}_{r\times r}(K)\)
jest dodatnio określona i odwracalna,
a rzut ortogonalny na \(U\) ma macierz
\[
 P_U=B(B^*B)^{-1}B^*.
\]
W szczególności
\[
 P_U^*=P_U,\quad
 P_U^2=P_U,\quad
 \operatorname{rk}P_U=r.
\]
Przy \(r=0\) należy rozumieć
\(P_U=0\) (puste iloczyny
macierzy mają właściwe rozmiary).
\end{twierdzenie}
\begin{proof}
Dla \(c\in K^r\) mamy
\[
 c^*Gc=c^*B^*Bc=\|Bc\|^2.
\]
Jeśli \(c\ne0\), niezależność
kolumn \(B\) oznacza \(Bc\ne0\),
więc \(c^*Gc>0\).
W szczególności gdy \(Gc=0\),
to \(c^*Gc=0\), a więc \(c=0\).
Macierz kwadratowa \(G\)
jest odwracalna.

Dla \(v\in K^n\) zapiszmy
\(c=G^{-1}B^*v\)
oraz \(u=Bc\in U\).
Wtedy
\[
 B^*(v-u)
 =B^*v-B^*BG^{-1}B^*v=0.
\]
Ponieważ kolumny \(B\)
rozpinają \(U\), równość
\(B^*(v-u)=0\) oznacza
\(v-u\in U^\perp\).
Na mocy jednoznaczności
rozkładu ortogonalnego
mamy \(u=P_Uv\),
co daje żądany wzór na macierz.

Ponieważ \(G^*=G\), także
\((G^{-1})^*=G^{-1}\).
Dlatego
\[
 P_U^*=B(G^{-1})^*B^*
      =BG^{-1}B^*=P_U.
\]
Ponadto
\[
 P_U^2
 =BG^{-1}(B^*B)G^{-1}B^*
 =BG^{-1}B^*=P_U.
\]
Jego obraz jest \(U\):
z jednej strony każdy obraz
leży w \(U\), a z drugiej
\(P_U(Bc)=Bc\) dla każdego \(c\).
Wobec tego \(\operatorname{rk}P_U=r\).
Dla \(r=0\) macierz \(B\)
ma zero kolumn i opisuje
podprzestrzeń zerową,
więc wzór oznacza rzut zerowy.
\end{proof}

\begin{przyklad}[Zespolony rzut obliczony macierzowo]
\label{prz:lin13-complex-projection}
Niech
\(a=(1,i)^{\mathsf T}\in\mathbb C^2\),
\(U=\operatorname{span}\{a\}\).
Dla macierzy \(B=a\) mamy
\(B^*B=2\), a więc
\[
 P_U=\frac12
 \begin{pmatrix}1&-i\\i&1\end{pmatrix}.
\]
Jest to macierz hermitowska
i idempotentna.
Dla \(v=(1,0)^{\mathsf T}\)
otrzymujemy
\[
 P_Uv=\tfrac12(1,i)^{\mathsf T},
 \qquad
 v-P_Uv=\tfrac12(1,-i)^{\mathsf T}.
\]
Iloczyn skalarny tych
dwóch wektorów jest zerowy:
\[
 \tfrac14(1+(-i)(-i))=0.
\]
W szczególności
\(\|v\|^2=1\),
\(\|P_Uv\|^2=\tfrac12\)
i \(\|v-P_Uv\|^2=\tfrac12\),
zgodnie ze wzorem Pitagorasa.
\end{przyklad}


\section{Równania normalne i metoda najmniejszych kwadratów}
\label{sec:lin13-least-squares}

\begin{definicja}[Zadanie najmniejszych kwadratów]
\label{def:lin13-ls}
Niech \(m,n\in\mathbb N_0\),
\(A\in\operatorname{Mat}_{m\times n}(K)\)
oraz \(y\in K^m\).
\emph{Zadaniem najmniejszych kwadratów}
jest znalezienie wszystkich \(x\in K^n\),
dla których wartość
\[
 \|Ax-y\|^2
 =\sum_{j=1}^m |(Ax)_j-y_j|^2
\]
jest najmniejsza.
Wektor \(Ax\) nazywamy
\emph{wektorem dopasowanym},
a \(r=y-Ax\) \emph{resztą}.
Nie zakładamy, że równanie
\(Ax=y\) jest rozwiązywalne.
\end{definicja}

\begin{twierdzenie}[Równania normalne]
\label{thm:lin13-normal-equations}
Niech \(A\in\operatorname{Mat}_{m\times n}(K)\)
oraz \(y\in K^m\), gdzie
\(m,n\in\mathbb N_0\).
Wtedy zadanie
\[
 \min_{x\in K^n}\|Ax-y\|^2
\]
ma przynajmniej jedno rozwiązanie.
Dla \(x_0\in K^n\)
następujące warunki są równoważne:
\begin{enumerate}[label=(\roman*)]
 \item \(x_0\) minimalizuje
       \(\|Ax-y\|^2\);
 \item reszta \(y-Ax_0\)
       należy do \((\operatorname{im}A)^\perp\);
 \item \(A^*(Ax_0-y)=0\);
 \item \(A^*Ax_0=A^*y\).
\end{enumerate}
Układ z (iv) nazywamy
\emph{układem równań normalnych}.
Wszystkie jego rozwiązania
mają ten sam wektor dopasowany:
\[
 Ax_0=P_{\operatorname{im}A}y.
\]
Jeśli \(x_0\) jest jednym rozwiązaniem,
to zbiorem wszystkich rozwiązań jest
\(x_0+\ker A\). W szczególności
rozwiązanie \(x_0\) jest jednoznaczne
dokładnie wtedy, gdy \(\ker A=\{0\}\).
\end{twierdzenie}
\begin{proof}
Obraz macierzy \(A\)
jest podprzestrzenią
skończonego wymiaru \(K^m\).
Z
Twierdzenia~\ref{thm:lin13-najlepsze}
wektor \(p=P_{\operatorname{im}A}y\)
jest jedynym wektorem w obrazie
minimalizującym \(\|p-y\|\).
Z definicji obrazu istnieje
\(x_0\in K^n\) takie, że
\(Ax_0=p\). To dowodzi
istnienia rozwiązania (i).

Dla dowolnego \(x\) mamy
\(Ax\in\operatorname{im}A\),
a na mocy wspomnianego
twierdzenia najlepsze
przybliżenie spośród takich
wektorów jest dokładnie \(p\).
Stąd (i) zachodzi wtedy i
tylko wtedy, gdy \(Ax_0=p\),
czyli reszta
\(y-Ax_0=y-p\)
jest ortogonalna do obrazu:
(i)\(\iff\)(ii).
Z
Twierdzenia~\ref{thm:lin13-kernel-image}
wynika
\((\operatorname{im}A)^\perp=\ker A^*\),
więc (ii)\(\iff\)(iii).
Równoważność (iii)\(\iff\)(iv)
wynika z liniowości operatora
sprzężonego.

Jeżeli \(x,x_0\) są rozwiązaniami,
to \(Ax=Ax_0=p\), więc
\(A(x-x_0)=0\) i
\(x-x_0\in\ker A\).
Odwrotnie, dodanie dowolnego
wektora z jądra do \(x_0\)
nie zmienia \(Ax_0\), więc
nie zmienia też reszty i
wartości minimalizowanej funkcji.
Rozwiązanie jest zatem
jednoznaczne wtedy i tylko wtedy,
gdy jądro jest trywialne.
Dla \(m=0\) obraz \(A\)
jest zerowy i wszystkie \(x\)
są rozwiązaniami; dla \(n=0\)
istnieje jedyny kandydat \(x=0\).
Dowód obejmuje oba przypadki.
\end{proof}

\begin{wniosek}[Przypadek kolumn niezależnych]
\label{cor:lin13-ls-full-rank}
Niech \(A\in\operatorname{Mat}_{m\times n}(K)\)
ma liniowo niezależne kolumny.
Wtedy układ równań normalnych
ma dokładnie jedno rozwiązanie,
dane wzorem
\[
 x_0=(A^*A)^{-1}A^*y.
\]
Wektor dopasowany i reszta
spełniają
\[
 Ax_0=A(A^*A)^{-1}A^*y,\qquad
 A^*(y-Ax_0)=0.
\]
\end{wniosek}
\begin{proof}
Niezależność kolumn oznacza
\(\ker A=\{0\}\).
Macierz \(A^*A\) jest
dodatnio określona:
dla każdego \(c\ne0\)
\[
 c^*A^*Ac=\|Ac\|^2>0.
\]
W szczególności jest odwracalna,
więc równania normalne mają
jedno rozwiązanie wyrażone
pierwszym wzorem.
Wstawiając je do \(Ax_0\)
otrzymujemy drugi wzór,
a trzeci to równania normalne.
Wszystkie te równości
pozostają prawdziwe dla \(n=0\)
przy standardowej konwencji
pustych macierzy.
\end{proof}

\begin{twierdzenie}[Rozwiązanie o najmniejszej normie]
\label{thm:lin13-ls-minnorm}
Niech \(A\in\operatorname{Mat}_{m\times n}(K)\),
\(y\in K^m\), gdzie \(m,n\ge0\).
Nawet gdy kolumny \(A\) są
liniowo zależne, istnieje
\emph{dokładnie jedno}
rozwiązanie najmniejszych kwadratów
\(x_{\min}\) o najmniejszej
normie euklidesowej.
Jest ono scharakteryzowane
dwoma warunkami
\[
 A^*(Ax_{\min}-y)=0,\qquad
 x_{\min}\in(\ker A)^\perp.
\]
Jeśli \(x_0\) jest dowolnym
rozwiązaniem najmniejszych kwadratów,
to
\[
 x_{\min}
 =P_{(\ker A)^\perp}x_0.
\]
\end{twierdzenie}
\begin{proof}
Na mocy
Twierdzenia~\ref{thm:lin13-normal-equations}
istnieje rozwiązanie \(x_0\);
wszystkie rozwiązania
stanowią zbiór \(x_0+\ker A\).
Przestrzeń \(K^n\) rozkłada się
ortogonalnie:
\[
 K^n=\ker A\oplus(\ker A)^\perp.
\]
Piszemy więc jednoznacznie
\(x_0=k_0+u_0\),
gdzie \(k_0\in\ker A\)
i \(u_0\in(\ker A)^\perp\).
Wektor \(u_0=x_0-k_0\)
również rozwiązuje zadanie,
bo \(Au_0=Ax_0\).
Każde inne rozwiązanie
ma postać \(u_0+k\)
dla pewnego \(k\in\ker A\).
Z ortogonalności
\[
 \|u_0+k\|^2=\|u_0\|^2+\|k\|^2.
\]
Norma jest zatem najmniejsza
dokładnie dla \(k=0\).
Wektor \(u_0\) jest jedynym
rozwiązaniem minimalnej normy
i należy do \((\ker A)^\perp\).
Pierwszy z warunków w tezie
jest równaniem normalnym,
a wzór na \(u_0\) wynika z
rzutu ortogonalnego na
\((\ker A)^\perp\).
\end{proof}

\begin{przyklad}[Dopasowanie prostej do trzech pomiarów]
\label{prz:lin13-regression}
Rozważmy rzeczywiste dane
\[
 (t_1,y_1)=(0,1),\qquad
 (t_2,y_2)=(1,2),\qquad
 (t_3,y_3)=(2,2).
\]
Szukamy współczynników \(a,b\)
prostej \(y=a+bt\),
minimalizując sumę
\(\sum_{j=1}^3(a+bt_j-y_j)^2\).
Dane odpowiadają macierzy i wektorowi
\[
 A=\begin{pmatrix}
 1&0\\1&1\\1&2
 \end{pmatrix},
 \qquad
 y=\begin{pmatrix}1\\2\\2\end{pmatrix}.
\]
Równania normalne to
\[
 \begin{pmatrix}3&3\\3&5\end{pmatrix}
 \begin{pmatrix}a\\b\end{pmatrix}
 =\begin{pmatrix}5\\6\end{pmatrix}.
\]
Wyznacznik macierzy
jest równy \(15-9=6>0\),
a jej odwrotność wynosi
\(\frac16\begin{pmatrix}5&-3\\-3&3\end{pmatrix}\).
Stąd
\[
 a=\frac76,\qquad b=\frac12.
\]
Wektor dopasowany i reszta to
\[
 Ax_0=
 \begin{pmatrix}7/6\\5/3\\13/6\end{pmatrix},
 \qquad
 r=y-Ax_0=
 \begin{pmatrix}-1/6\\1/3\\-1/6\end{pmatrix}.
\]
Bezpośrednio sprawdzamy
\[
 A^{\mathsf T}r
 =\begin{pmatrix}
 -1/6+1/3-1/6\\
 0+1/3-2/6
 \end{pmatrix}
 =\begin{pmatrix}0\\0\end{pmatrix}.
\]
Reszta jest więc ortogonalna
do każdej kolumny \(A\).
Minimalna suma kwadratów
błędów wynosi
\[
 \|r\|^2
 =\frac1{36}+\frac19+\frac1{36}
 =\frac16.
\]
W tym przykładzie
układ \(Ax=y\) nie ma
rozwiązania dokładnego,
mimo że zadanie najmniejszych
kwadratów ma jednoznaczne
rozwiązanie \(x_0\).
\end{przyklad}

\begin{przyklad}[Macierz niepełnego rzędu]
\label{prz:lin13-rank-deficient}
Niech
\[
 A=\begin{pmatrix}1&1\\2&2\end{pmatrix},
 \qquad
 y=\begin{pmatrix}1\\0\end{pmatrix}.
\]
Obraz \(A\) to prosta
rozpięta przez \(a=(1,2)^{\mathsf T}\).
Rzut \(y\) na tę prostą wynosi
\[
 P_{\operatorname{im}A}y
 =\frac{a^{\mathsf T}y}{a^{\mathsf T}a}\,a
 =\frac15\begin{pmatrix}1\\2\end{pmatrix}.
\]
Wszystkie rozwiązania
najmniejszych kwadratów
spełniają więc warunek
\(x_1+x_2=1/5\).
Macierz \(A^{\mathsf T}A
=\begin{pmatrix}5&5\\5&5\end{pmatrix}\)
jest osobliwa, toteż nie wolno
stosować do niej wzoru
z odwrotnością.
Spośród rozwiązań
\(x=(t,\frac15-t)\)
najmniejszą normę ma
\[
 x_{\min}=\begin{pmatrix}1/10\\1/10\end{pmatrix},
\]
gdyż jest prostopadły do
\(\ker A=\operatorname{span}\{(1,-1)\}\).
Zgadza się to z
Twierdzeniem~\ref{thm:lin13-ls-minnorm}.
\end{przyklad}

\section{Złożenia rzutów ortogonalnych}
\label{sec:lin13-rzuty-zlozenia}

\begin{twierdzenie}[Kiedy iloczyn rzutów jest rzutem ortogonalnym]
\label{thm:lin13-commuting-projections}
Niech \((V,h)\) będzie skończenie
wymiarową przestrzenią
euklidesową lub unitarną,
a \(P=P_U\), \(Q=P_W\)
rzutami ortogonalnymi na
podprzestrzenie \(U,W\subseteq V\).
Następujące warunki są równoważne:
\begin{enumerate}[label=(\roman*)]
 \item \(PQ=QP\);
 \item \(PQ\) jest operatorem
       sprzężonym do samego siebie;
 \item \(PQ\) jest rzutem
       ortogonalnym.
\end{enumerate}
Jeśli te warunki zachodzą, to
\[
 PQ=P_{U\cap W}.
\]
Nie zakładamy, że \(U\) i \(W\)
są do siebie ortogonalne.
\end{twierdzenie}
\begin{proof}
Ponieważ \(P^*=P\) i
\(Q^*=Q\), z
Stwierdzenia~\ref{prop:lin13-adjoint-kalkul}
wynika
\[
 (PQ)^*=Q^*P^*=QP.
\]
Zatem (i)\(\iff\)(ii).
Jeśli \(P,Q\) przemieniają się,
to
\[
 (PQ)^2=PQPQ=P^2Q^2=PQ
\]
i \((PQ)^*=PQ\).
Na mocy
Twierdzenia~\ref{thm:lin13-projection-characterization}
operator \(PQ\) jest więc
rzutem ortogonalnym,
co dowodzi (i)\(\Rightarrow\)(iii).
Z kolei każdy rzut ortogonalny
jest sprzężony do samego
siebie, więc (iii)\(\Rightarrow\)(ii).

Przy przemienności \(PQ=QP\)
obraz \(PQ\) leży w \(U\),
gdyż \(P(PQv)=P^2Qv=PQv\).
Leży też w \(W\), bo
\(Q(PQv)=P Q^2v=PQv\).
Zatem
\(\operatorname{im}(PQ)\subseteq U\cap W\).
Jeśli \(x\in U\cap W\), to
\(Px=Qx=x\), więc \(PQx=x\).
Stąd \(U\cap W\subseteq
\operatorname{im}(PQ)\).
Obraz \(PQ\) jest dokładnie
przecięciem, więc \(PQ=P_{U\cap W}\).
\end{proof}

\begin{przyklad}[Dwa rzuty, które nie przemieniają się]
\label{prz:lin13-noncommuting}
W \(\mathbb R^2\) rozważmy
\[
 U=\operatorname{span}\{(1,0)\},
 \qquad
 W=\operatorname{span}\{(1,1)\}.
\]
Ich rzuty ortogonalne mają macierze
\[
 P=\begin{pmatrix}1&0\\0&0\end{pmatrix},
 \qquad
 Q=\frac12\begin{pmatrix}1&1\\1&1\end{pmatrix}.
\]
Bezpośrednie mnożenie daje
\[
 PQ=\frac12\begin{pmatrix}1&1\\0&0\end{pmatrix},
 \qquad
 QP=\frac12\begin{pmatrix}1&0\\1&0\end{pmatrix}.
\]
Zatem \(PQ\ne QP\).
Ponieważ \(PQ\) nie jest
macierzą symetryczną,
nie jest też rzutem
ortogonalnym.
W tym przypadku
\(U\cap W=\{0\}\),
a więc \(P_{U\cap W}=0\ne PQ\).
\end{przyklad}

\begin{stwierdzenie}[Dopełnienie przecięcia w skończonym wymiarze]
\label{prop:lin13-perp-intersection}
Niech \((V,h)\) będzie
skończenie wymiarową przestrzenią
euklidesową lub unitarną,
a \(U,W\subseteq V\)
podprzestrzeniami.
Wtedy
\[
 (U\cap W)^\perp=U^\perp+W^\perp.
\]
Równość jest rozumiana jako
algebraiczna suma podprzestrzeni.
\end{stwierdzenie}
\begin{proof}
Z
Stwierdzenia~\ref{prop:lin13-dopelnienia}
mamy
\[
 (U^\perp+W^\perp)^\perp
 =(U^\perp)^\perp\cap(W^\perp)^\perp.
\]
W skończonym wymiarze
Twierdzenie~\ref{thm:lin13-rozklad}
daje
\((U^\perp)^\perp=U\)
oraz
\((W^\perp)^\perp=W\).
Wynika stąd
\((U^\perp+W^\perp)^\perp=U\cap W\).
Ponowne dopełnienie
ortogonalne, zastosowane
do lewej i prawej strony,
daje tezę, ponieważ w tej
przestrzeni każda podprzestrzeń
jest równa swojemu
podwójnemu dopełnieniu.
\end{proof}

\section{Podsumowanie}
\label{sec:lin13-podsumowanie}

Dla skończenie wymiarowej
podprzestrzeni \(U\)
w dowolnej przestrzeni
euklidesowej lub unitarnej
istnieje jednoznaczny
rozkład \(V=U\oplus U^\perp\).
Rzut \(P_U\) wybiera pierwszy
składnik tego rozkładu,
jest idempotentny i zachowuje
właściwą symetrię iloczynu
skalarnego. Rzut jest
najlepszym przybliżeniem
wektora przez elementy \(U\),
a reszta jest do \(U\)
ortogonalna.
W skończonym wymiarze
operator sprzężony \(A^*\)
wyznacza związki między
obrazem a jądrem:
\((\operatorname{im}A)^\perp=\ker A^*\).
Rzut ortogonalny jest dokładnie
operatorem \(P\), dla którego
\(P^2=P=P^*\).
Dla przestrzeni rozpiętej przez
niezależne kolumny macierzy \(B\)
otrzymujemy obliczalną
postać \(P=B(B^*B)^{-1}B^*\).
Równania normalne
\(A^*Ax=A^*y\) charakteryzują
wszystkie rozwiązania
najmniejszych kwadratów;
nawet dla macierzy niepełnego
rzędu istnieje jednoznaczne
rozwiązanie o najmniejszej
normie, choć sam układ
może mieć wiele rozwiązań.
Iloczyn dwóch rzutów
ortogonalnych jest rzutem
ortogonalnym dokładnie wtedy,
gdy rzuty się przemieniają.
Te wyniki stanowią podstawę
dalszej teorii operatorów,
rozkładów macierzy i metod
numerycznych.
